\documentclass[../../main.tex]{subfiles}
\begin{document}
\begin{oframed}
  \subfile{problem}
\end{oframed}

{\bf [解答]}

\paragraph{(i)}

$f(x)=(1+x)^{\alpha}-\left(1+\dfrac{\alpha}{2}x\right)$とおく．
\begin{align}
  f'(x)=\alpha\left[\left(\dfrac{1}{1+x}\right)^{1-\alpha}-\dfrac{1}{2}\right]
\end{align}
である．
\begin{align}
  &0\le x\le1 &0<\alpha<1
\end{align}
から，$\cfrac{1}{2}\le\cfrac{1}{1+x}\le1$，$0<1-\alpha<1$であることより，
\begin{align}
  \dfrac{1}{2}<\left(\dfrac{1}{2}\right)^{1-\alpha}\le\left(\dfrac{1}{1+x}\right)^{1-\alpha}
\end{align}
従って$f'(x)>0$だから，$f(x)$は単調増加．これと$f(0)=0$から，区間内で$f(x)\ge0$である．故に示された．$\Box$

\paragraph{(ii)}
$2^{200}=(2^{10})^{20}$だから，$2^{10}>10^3$ゆえ，
\begin{align}
  2^{200}>10^{60}\label{eq:1}
\end{align}
である．

また，(1)で$(x,\alpha)=(1,1/20)$とすることにより，(これらは条件を満たす．)
\begin{align}
  &2^{1/20}\ge 1+\dfrac{1}{40}>1+\dfrac{24}{1000} \nonumber\\
  \Longleftrightarrow &2^{10}<2^{1/20}10^3 \nonumber\\
  \therefore \ &2^{200}<2×10^{60} \label{eq:2}
\end{align}
を得る．\cref{eq:1}，\cref{eq:2}より，
\begin{align}
  10^{60}<2^{200}<2×10^{60}\label{eq:3}
\end{align}
であるから，$61$桁で，最高位は$1$である．$\cdots$(答)

\paragraph{(iii)}
\cref{eq:3}の各辺正より，常用対数を取る.以下$\log_{10}2=a$とする．
\begin{align}
  &60<200a<60+a  \\
  \Longleftrightarrow &\dfrac{3}{10}<a<\dfrac{60}{199} \\
  \therefore \ &0.300<a<0.3015\cdots
\end{align}
であるから，題意は示された．$\Box$
\newpage
\end{document}
